% ==================== 附录 ====================
\clearpage
\section*{附\quad 录}
\addcontentsline{toc}{section}{附\quad 录}

\subsection*{附录 A：核心数据类定义（节选）}
\addcontentsline{toc}{subsection}{附录 A：核心数据类定义（节选）}

本附录节选自 \code{classmind/models.py}，用于说明业务对象在代码层的
承载方式。完整源码请参见项目仓库 \code{classmind/models.py}。

\begin{lstlisting}[style=cmplain]
# classmind/models.py 节选
from __future__ import annotations
from dataclasses import asdict, dataclass, field
from typing import Any


@dataclass(frozen=True)
class TimeSlot:
    id: str
    day: str
    period: str
    order: int


@dataclass(frozen=True)
class Teacher:
    id: str
    name: str
    qualifications: tuple = ()
    unavailable: tuple = ()
    preferred_slots: tuple = ()


@dataclass(frozen=True)
class Room:
    id: str
    name: str
    campus: str
    capacity: int
    equipment: tuple = ()
    unavailable: tuple = ()


@dataclass(frozen=True)
class ClassGroup:
    id: str
    name: str
    size: int
    unavailable: tuple = ()
    preferred_slots: tuple = ()


@dataclass(frozen=True)
class CourseRequest:
    id: str
    name: str
    subject: str
    class_id: str
    sessions: int
    qualified_teacher_ids: tuple = ()
    required_equipment: tuple = ()
    allowed_slots: tuple = ()


@dataclass(frozen=True)
class ScheduledLesson:
    lesson_id: str
    course_id: str
    course_name: str
    class_id: str
    class_name: str
    teacher_id: str
    teacher_name: str
    room_id: str
    room_name: str
    slot_id: str
    day: str
    period: str
\end{lstlisting}

\subsection*{附录 B：演示数据示例（节选）}
\addcontentsline{toc}{subsection}{附录 B：演示数据示例（节选）}

以下片段来自 \code{data/demo.json}，仅保留字段结构与一两条记录，
完整数据请参见项目仓库 \code{data/demo.json}。

\begin{lstlisting}[style=cmjson]
// data/demo.json 节选
{
  "time_slots": [
    {"id":"MON_0800","day":"周一","period":"08:00-09:30","order":1},
    {"id":"FRI_1530","day":"周五","period":"15:30-17:00","order":20}
  ],
  "teachers": [
    {"id":"T01","name":"张老师","qualifications":["数学"],
     "unavailable":["WED_1000"],
     "preferred_slots":["MON_0800","TUE_0800","THU_1000"]},
    {"id":"T02","name":"李老师","qualifications":["数学","物理"],
     "unavailable":["TUE_1000"]}
  ],
  "rooms": [
    {"id":"R101","name":"思源教室","campus":"鼓楼校区",
     "capacity":30,"equipment":["投影","白板"]}
  ],
  "classes": [
    {"id":"C01","name":"高一数学A班","size":28}
  ],
  "courses": [
    {"id":"CR01","name":"高一数学提升","subject":"数学",
     "class_id":"C01","sessions":2,
     "qualified_teacher_ids":["T01","T02"],
     "required_equipment":["投影"]}
  ]
}
\end{lstlisting}

\subsection*{附录 C：API 接口详细定义}
\addcontentsline{toc}{subsection}{附录 C：API 接口详细定义}

本节列出当前原型已实现的核心接口，字段定义以服务实现为准。

\subsubsection*{C.1 GET /api/health}

\textbf{用途：}健康检查。\textbf{请求：}无。

\begin{lstlisting}[style=cmjson]
{
  "status": "ok",
  "service": "ClassMind",
  "solver": "OR-Tools CP-SAT"
}
\end{lstlisting}

\subsubsection*{C.2 GET /api/plans}

\textbf{用途：}获取多套候选方案。\textbf{请求：}无。
\textbf{响应：}包含方案列表，每条方案附 \code{status}、\code{metrics}、
\code{conflicts}、\code{schedule}、\code{scorecard} 字段。

\subsubsection*{C.3 GET /api/dashboard}

\textbf{用途：}获取决策驾驶舱所需数据。
\textbf{响应：}单个 \code{SolveResult}，额外含 \code{teacher\_load} 字段。

\subsubsection*{C.4 POST /api/solve}

\textbf{用途：}对自定义业务数据求解。
\textbf{请求体：}与 \code{data/demo.json} 同结构，可选 \code{strategy} 字段。
\textbf{响应：}单个 \code{SolveResult}。

\subsubsection*{C.5 POST /api/reschedule}

\textbf{用途：}教师请假最小影响调课。

\begin{lstlisting}[style=cmjson]
{ "teacher_id": "T01", "slot_id": "MON_0800" }
\end{lstlisting}

\textbf{响应：}包含 \code{before}、\code{after}、\code{diff}、\code{impact}
四个字段，分别表示调课前后方案、变更明细及影响对象。

\subsection*{附录 D：核心算法节选}
\addcontentsline{toc}{subsection}{附录 D：核心算法节选}

完整源码托管在 GitHub：
\begin{center}
  \code{github.com/Lenny-lab/classmind-scheduler}
\end{center}

本节给出 \code{solver.py} 求解入口与权重表片段，其余模块
（\code{api.py / service.py / validator.py / io.py / models.py / cli.py}）
请参见上述仓库的 \code{classmind/} 目录。

\begin{lstlisting}[style=cmplain]
# classmind/solver.py 节选
from ortools.sat.python import cp_model


STRATEGY_WEIGHTS = {
    "student":    {"preference": 10, "stability": 2,
                   "efficiency": 1, "change": 0, "compactness": 3},
    "teacher":    {"preference": 2,  "stability": 10,
                   "efficiency": 1, "change": 0, "compactness": 3},
    "efficiency": {"preference": 1,  "stability": 2,
                   "efficiency": 10,"change": 0, "compactness": 3},
    "balanced":   {"preference": 5,  "stability": 5,
                   "efficiency": 5, "change": 0, "compactness": 3},
    "reschedule": {"preference": 1,  "stability": 1,
                   "efficiency": 1, "change": 100, "compactness": 0},
}


def solve_schedule(problem, time_limit_seconds=15.0,
                   strategy="balanced", baseline=None):
    # 1. 业务数据校验
    data_errors = validate_problem(problem)
    if data_errors:
        return SolveResult(status="INVALID_DATA", conflicts=data_errors)

    # 2. 策略与权重选择
    if strategy not in STRATEGY_WEIGHTS:
        return SolveResult(status="INVALID_DATA",
                           conflicts=[f"未知策略: {strategy}"])
    weights = STRATEGY_WEIGHTS[strategy]

    # 3. 变量预筛选与硬约束建模（仅保留可行组合）
    model = cp_model.CpModel()
    variables = {}
    lesson_vars = defaultdict(list)
    for lesson_id, course in lesson_specs:
        for (teacher_id, room_id, slot_id) in feasible_combos(course):
            key = (lesson_id, teacher_id, room_id, slot_id)
            variables[key] = model.new_bool_var("x_" + "_".join(key))
            lesson_vars[lesson_id].append(variables[key])

    # 4. 互斥与"恰好一次"约束
    for lesson_id, _ in lesson_specs:
        model.add_exactly_one(lesson_vars[lesson_id])
    for resource in ("teacher", "room", "class"):
        for slot_id in slot_ids:
            model.add_at_most_one(...)

    # 5. 加权软目标
    objective_terms = []
    objective_terms += _per_lesson_soft_costs(...)
    objective_terms += _stability_cost(...)
    model.minimize(sum(objective_terms))

    # 6. 求解 + 独立校验
    solver = cp_model.CpSolver()
    solver.parameters.max_time_in_seconds = time_limit_seconds
    status_code = solver.solve(model)
    schedule = extract_schedule(solver, variables, ...)
    conflicts = validate_schedule(problem, schedule)
    return SolveResult(status=..., schedule=schedule,
                       conflicts=conflicts, metrics={...})
\end{lstlisting}

\subsection*{附录 E：术语表}
\addcontentsline{toc}{subsection}{附录 E：术语表}

\begin{table}[h]
  \centering
  \caption{术语表}
  \label{tab:glossary}
  \small
  \begin{tabularx}{\textwidth}{L{3.2cm} Y}
    \toprule
    \textbf{术语} & \textbf{说明} \\
    \midrule
    CP-SAT
      & Constraint Programming-SAT，Google OR-Tools 中的混合求解引擎，
        结合约束传播、SAT 冲突学习与线性规划松弛。\\
    CP & Constraint Programming，约束规划。\\
    CSP & Constraint Satisfaction Problem，约束满足问题。\\
    SAT & Boolean Satisfiability Problem，布尔可满足性问题。\\
    OR-Tools & Google 开源的运筹学工具包，Apache 2.0 协议。\\
    硬约束
      & 必须严格满足的约束，违反则方案无效（如教师同时段冲突数必须为 0）。\\
    软目标
      & 可以被违反但应尽量满足的目标（如学生偏好尽量命中）。\\
    策略权重
      & 不同维度的软目标在总目标中的相对重要性（0--10）。\\
    策略
      & 一组策略权重的命名（balanced / student / teacher / efficiency / reschedule）。\\
    稳定性绝对偏差
      & 教师日均课次相对目标值的绝对偏差之和，越低表示日负载越均匀。\\
    调课
      & 在原方案基础上做最小影响的修改。\\
    决策变量
      & 求解器可以控制的变量（如课次到时间槽的分配）。\\
    预筛选
      & 在生成决策变量前就剔除明显不可行的组合。\\
    辅助变量
      & 求解器引入的额外变量，用于表达复杂约束。\\
    教务老师
      & K12 培训机构中负责排课的老师。\\
    驾驶舱
      & Dashboard，集中展示关键指标的页面。\\
    \bottomrule
  \end{tabularx}
\end{table}
